\chapter{Verificación formal de identidades algebraicas finitas}
\label{app:lean}

Las igualdades finitas siguientes se formalizan en Lean y se deciden por
reducción del núcleo lógico, sin axiomas externos adicionales. Los teoremas
demuestran:
\begin{longtable}{p{.40\textwidth}p{.50\textwidth}}
\toprule
Teorema & Contenido\\
\midrule
\endhead
\texttt{aw\_symmetric}
& simetría de la matriz Paley--Witt \(A_W\);\\
\texttt{aw\_square\_minus\_identity}
& \(A_W^2=-I_6\) sobre \(\mathbb F_3\);\\
\texttt{x\_sq\_add\_one\_has\_no\_root}
& irreducibilidad de \(x^2+1\) sobre \(\mathbb F_3\);\\
\texttt{dual\_has\_nonzero\_nilpotent}
& existencia de un nilpotente en el álgebra parabólica;\\
\texttt{split\_has\_nontrivial\_idempotent}
& existencia de un idempotente no trivial en el álgebra hiperbólica;\\
\texttt{canonical\_charge}
& \(\sum_{m=1}^{12}K_m=6263\);\\
\texttt{hadamard\_orbit\_one/two/three}
& las tres órbitas de Hadamard producen el estado dodecafásico publicado.\\
\bottomrule
\end{longtable}

La formalización recibe las matrices y el vector dodecafásico como datos
tipados y verifica sus consecuencias algebraicas. La construcción del estado
firmado y la enumeración de (104\,976) semillas pertenecen a los teoremas
combinatorios previos y no se sustituyen por esta reducción formal.
